% Generated draft: requires mathematical and source review.
\documentclass[11pt]{article}
\usepackage[margin=25mm]{geometry}
\usepackage{fontspec}
\setmainfont{Latin Modern Roman}
\usepackage{amsmath,amssymb,mathtools,unicode-math}
\setmathfont{Latin Modern Math}
\usepackage{enumitem,needspace,xurl,hyperref}
\setlength{\parindent}{0pt}
\setlength{\parskip}{5pt}
\setlength{\emergencystretch}{3em}
\newcommand{\N}{\mathbb{N}}
\newcommand{\Z}{\mathbb{Z}}
\newcommand{\Q}{\mathbb{Q}}
\newcommand{\R}{\mathbb{R}}
\newcommand{\C}{\mathbb{C}}
\newcommand{\F}{\mathbb{F}}
\newcommand{\A}{\mathbb{A}}
\newcommand{\PP}{\mathbb{P}}
\newcommand{\E}{\mathbb{E}}
\newcommand{\eps}{\varepsilon}
\newcommand{\dd}{\,\mathrm{d}}
\begin{document}
\section*{UnsolvedMath ID: 30001087}
\textbf{Source problem code:} catalogue entry
\subsection{Definitions and mathematical statement}
Let \(f\in\C[z]\) have \(f(0)=0\) and \(f'(0)=\lambda=e^{2\pi i\theta}\), where \(\theta\in\R\setminus\Q\). The fixed point is analytically linearizable if there is a local biholomorphic map \(h\), with \(h(0)=0\) and \(h'(0)=1\), satisfying
\[
 f(h(w))=h(\lambda w)
\]
near zero. When it is linearizable, let \(\Delta_f\) be the Siegel disk containing zero, namely the maximal corresponding rotation domain. Its conformal radius is
\[
 r(f)=|\psi'(0)|,
 \qquad\psi:\mathbb D\xrightarrow{\cong}\Delta_f,\quad\psi(0)=0.
\]
This is independent of the choice of \(\psi\). It is not by definition the Euclidean distance from zero to \(\partial\Delta_f\), and no radius is assigned here when there is no Siegel disk.

Let \(q_j\) be the denominators of the continued-fraction convergents of \(\theta\) modulo integers. The Brjuno condition is
\[
 B(\theta)=\sum_{j\geq0}\frac{\log q_{j+1}}{q_j}<\infty.
\]
The broad requests are to characterize linearizability for the polynomial and determine or estimate \(r(f)\) when defined. The precise conjecture attributed to Douady is
\[
 \forall d\geq3\ \forall f\in\C[z]\text{ with }\deg f=d,\ f(0)=0,
 \ f'(0)=e^{2\pi i\theta},\ \theta\notin\Q,
\]
\[
 f\text{ linearizable at }0\quad\Longrightarrow\quad B(\theta)<\infty.
\]
This is necessity for each individual polynomial, not just a statement that the multiplier condition guarantees linearization simultaneously for all germs or all polynomials. No restriction to saturated, unicritical, or other special polynomial classes is present, and degree two is outside the conjecture's stated higher-degree range.

The conformal-radius request supplies no quantitative target or coefficient normalization. It cannot in general be read as a radius determined by \(\theta\) alone: for \(a\ne0\), the polynomial \(f_a(z)=a^{-1}f(az)\) has the same multiplier and its Siegel-disk radius scales as \(r(f_a)=r(f)/|a|\). Thus size estimates must retain the actual polynomial or a specified normalization.

The supplied assessment dated 2026-08-21 reports the general cubic and higher-degree necessity question as open, with substantial special cases, and lists saturated-polynomial and cubic-Siegel references. It supplies no complete special-case hypotheses or theorem resolving the full displayed assertion.
\subsection{Short English statement}
Douady's conjecture says that any linearizable irrationally indifferent fixed point of a polynomial of degree at least three must have Brjuno rotation number. The record also asks for conformal-radius information; such a radius depends on the polynomial's scale and is only defined when its Siegel disk exists.
\subsection{Sources}
\begin{itemize}
\item Local UnsolvedMath registry, permanent ID \(30001087\), exact source code \texttt{catalogue entry}, proposed year 2008. The original explicitly attributes the necessity conjecture to Douady; no separate proposer field is supplied.
\item \emph{Trends and Developments in Complex Dynamics (2009)}, Oberwolfach Reports, DOI 10.4171/owr/2008/48, pp. 2703--2762. \url{https://doi.org/10.4171/owr/2008/48}.
\item Ch\'eritat, ``linearizability of saturated polynomials,'' \url{https://arxiv.org/abs/1507.02666}.
\item Supplied reference labelled ``Cubic Siegel polynomials, current status'': \url{https://link.springer.com/article/10.1007/s12220-026-02376-2}.
\item The supplied \texttt{corrected\_verified} note reports preservation of the full statement with notation cleanup. The open label and assessment checked 2026-08-21 are supplied dated curation, not independently verified; exact special cases and current status require review.
\end{itemize}
\end{document}
